% TOP_PROBLEM: {"id": 3057, "title": "Perfect 2-error-correcting codes over arbitrary finite alphabets.", "category_id": 2, "category": {"id": 2, "name": "combinatorics", "display_name": "Combinatorics", "description": "Counting problems, graph theory, discrete structures.", "slug": "combinatorics", "order_index": 2, "created_at": "2026-07-31T15:26:25.671Z"}, "status": "open", "problem_number": "OPG-37196", "set_id": 12}
% FIRST_POSTED: 2026-10-01
% ENTRY_KIND: statement_only
% UnsolvedMath ID 3057; source catalogue code OPG-37196
% !TeX program = lualatex
% Definition and sources only; this file is not an LLM solution attempt.
\documentclass[11pt,a4paper]{article}
\usepackage[margin=25mm]{geometry}
\usepackage{fontspec}
\setmainfont{lmroman10-regular.otf}[BoldFont=lmroman10-bold.otf,
 ItalicFont=lmroman10-italic.otf,BoldItalicFont=lmroman10-bolditalic.otf]
\usepackage{amsmath,amssymb,mathtools}
\usepackage{unicode-math}
\setmathfont{latinmodern-math.otf}
\AtBeginDocument{\let\setminus\smallsetminus}
\usepackage{xurl,hyperref}
\hypersetup{colorlinks=true,urlcolor=blue}
\setcounter{secnumdepth}{0}
\setlength{\parindent}{0pt}
\setlength{\parskip}{5pt}
\setlength{\emergencystretch}{3em}
\begin{document}
\section*{Perfect 2-error-correcting codes over arbitrary finite alphabets.}\label{3057}
{\small UnsolvedMath ID 3057 · OPG-37196 · Combinatorics · OpenGarden}
\subsection{Definitions and mathematical statement}
Let $Q$ be a finite alphabet of cardinality $q\ge2$ and let $n$ be a positive integer. A code is a subset $C\subseteq Q^n$; linearity, a field structure on $Q$, and prime-power alphabet size are not assumed. The Hamming distance $d(x,y)$ counts coordinates in which two words differ. Write $B_2(c)=\{x\in Q^n:d(x,c)\le2\}$. A perfect two-error-correcting code is a code for which these balls, as $c$ ranges over $C$, partition $Q^n$: every word has exactly one codeword at distance at most two.

The nontrivial classification problem asks whether every such code with $|C|>2$ is equivalent to the ternary Golay code of length $11$. Equivalence permits a permutation of coordinates and, independently in each coordinate, a permutation of the alphabet. The ternary Golay code is the linear $[11,6,5]$ code over $\mathbb F_3$, namely a six-dimensional subspace with minimum nonzero Hamming weight five; the standard code of these parameters is unique up to this equivalence.

In particular, the predicted parameters are $q=3$, $n=11$, and $|C|=3^6$. A necessary counting identity is
\[
 |C|\left(1+n(q-1)+\binom n2(q-1)^2\right)=q^n.
\]
Satisfying this identity alone does not produce a code: unique decoding requires the actual disjoint covering property. The question ranges over arbitrary finite alphabets and arbitrary, possibly nonlinear, codes.
\subsection{Short English statement}
Must every nontrivial perfect code correcting exactly two Hamming errors be a ternary Golay code?
\subsection{Sources}
\begin{itemize}
\item[S1] ULAM.AI and the UnsolvedMath Contributors, supplied problems.json, record 3057 (OPG-37196), OpenGarden. Catalogue identity and source transcription; CC BY 4.0.
\url{https://huggingface.co/datasets/ulamai/UnsolvedMath}.
\item[S2] Open Problem Garden, Perfect 2-error-correcting codes over arbitrary finite alphabets.. Original problem and discussion.
\url{http://www.openproblemgarden.org/op/perfect_2_error_correcting_codes_over_arbitrary_finite_alphabets}.
\item[S3] M. A. Bennett, Perfect 2-codes over arbitrary alphabets (2026).
\url{https://arxiv.org/abs/2607.27555}.
\end{itemize}
\end{document}
